\section{The compact avoiding set}\label{sec:global}

The periodic hitting sets cover scales in $[1,2]$. We now combine their
dyadic dilations, with summable densities, to cover every nonzero scale
while removing arbitrarily little measure from $[0,1]$.

\begin{proof}[Proof of Theorem~\ref{thm:main}]
Fix $\eta\in(0,1)$ and put
\[
 p_j=\frac{\eta}{24}\,2^{-j}\qquad(j=1,2,\ldots).
\]
Each $p_j$ belongs to $(0,1)$. By Lemma~\ref{lem:periodic}, choose an
open $1$-periodic set $H_j\subset\R$ with $\rho(H_j)\le6p_j$ such that
\begin{equation}\label{global:normalized-hits}
 \forall y\in\R\ \forall t\in[1,2]\ \exists n\in\mathbb Z_{\ge1}:
 \qquad y+t2^{-n}\in H_j.
\end{equation}
All these sets are fixed before any center or scale is chosen. Define
\[
 C=\bigcup_{j\ge1}\bigl(2^{-j}H_j\cup(-2^{-j}H_j)\bigr),
 \qquad E=[0,1]\setminus C.
\]
Dilations and reflections preserve openness, so $C$ is open. Its paired
definition gives $-C=C$. Thus $E$ is a closed subset of $[0,1]$, and is
compact.

For each $j$, periodicity and a change of variables give the exact
identities
\begin{align*}
 m\bigl((2^{-j}H_j)\cap[0,1]\bigr)
 &=2^{-j}m\bigl(H_j\cap[0,2^j]\bigr)=\rho(H_j),\\
 m\bigl((-2^{-j}H_j)\cap[0,1]\bigr)
 &=2^{-j}m\bigl(H_j\cap[-2^j,0]\bigr)=\rho(H_j).
\end{align*}
Indeed, each interval on the right contains exactly $2^j$ unit periods;
their endpoints do not affect Lebesgue measure. Countable subadditivity
therefore yields
\[
 m(C\cap[0,1])\le2\sum_{j\ge1}\rho(H_j)
     \le12\sum_{j\ge1}p_j=\frac{\eta}{2}.
\]
Consequently
\[
 m(E)=1-m(C\cap[0,1])\ge1-\frac{\eta}{2}>1-\eta.
\]

It remains to show that every signed affine copy of $D$ meets $C$.
Let $x\in\R$ and first suppose that $s>0$. Choose an integer $k$ with
\[
 t=2^k s\in[1,2),
\]
and set $j=\max\{1,k\}$. Apply~\eqref{global:normalized-hits} to the
center $2^j x$ and the normalized scale $t$. There is an integer $n\ge1$
such that
\[
 2^j x+t2^{-n}\in H_j.
\]
Multiplying by $2^{-j}$ gives
\[
 x+s2^{-(n+j-k)}\in2^{-j}H_j\subset C.
\]
The index $N=n+j-k$ is an integer and satisfies $N\ge n\ge1$, because
$j\ge k$. Hence this point belongs to the original copy $x+sD$.
The argument applies to every real center and every positive scale,
including arbitrarily small scales, for which larger values of $j$ are
available.

If $s<0$, apply the positive-scale conclusion to the center $-x$ and
scale $-s$. For some integer $N\ge1$,
\[
 -(x+s2^{-N})=(-x)+(-s)2^{-N}\in C.
\]
Since $-C=C$, it follows that $x+s2^{-N}\in C$ as well. In either case,
the copy $x+sD$ has a point in $C$, which is disjoint from $E$.
Thus the fixed compact set $E$ excludes every such copy, as required.
\end{proof}
